\section{Generalization to Other Portable Devices}\label{sec:general}
The equations of \cref{sec:model} hold for any regulated load on a lithium-ion pack; only constants change. For
$N_s$ cells in series and $N_p$ in parallel of capacity $Q_{\mathrm{cell}}$ (pack $Q=N_pQ_{\mathrm{cell}}$, mean OCV
$\bar U$ per cell), the nominal energy is $E_n=N_sQ\bar U$ and $\TTE\approx E_u/\bar P$ with $E_u$ slightly below
$E_n$; $R\propto(N_s/N_p)/Q_{\mathrm{cell}}$ (RP-63: $R_0Q=0.28\,\Omega$\,A\,h), $C_{\mathrm{th}}\propto$ mass,
$R_{\mathrm{th}},R_{\mathrm{skin}}\propto1/(h_cA_{\mathrm{surf}})$ (heat-transfer coefficient $h_c$), so the
sustainable power is $P_{\mathrm{sust}}\propto(T_{\lim}-\Tamb)\,h_cA_{\mathrm{surf}}$. Display laws differ
(OLED $\propto A\,g(b)\,\mathrm{APL}$, LCD $P_{\mathrm{panel}}+c_1b+c_2b^2$, e-paper an energy per page); usage
needs its own activity set but keeps F6--F8.

Five dimensionless groups tell which model layer a device needs (\cref{tab:general}): ohmic loss
$\Pi_1=PR_\Sigma/\bar U^2$ (the first term of F3), throttling $\Pi_2=P_{\mathrm{dem}}/P_{\mathrm{sust}}$ (F5),
thermal memory $\Pi_3=\tau_{\mathrm{th}}/\TTE$ (thermal time constant $\tau_{\mathrm{th}}$), burst margin
$\Pi_4=\Delta P_b/\Delta P_{b,\mathrm{safe}}$ (F4; brown-out at 1) and usage spread $\Pi_5=\mathrm{CV}_{\TTE}$ of F7,
$\approx\sqrt{2\tau_c/\mu}\,\mathrm{CV}_P$ ($\tau_c$: correlation time of the power, $\mu$: mean TTE). Per cell,
$\Pi_1=c\,R_\Sigma Q/\bar U$ and $\Pi_4$ depend only on the C-rate $c$ and $R_\Sigma Q$, not on pack size.

\begin{table}[t]\centering\small
\caption{Dimensionless groups. RP-63: $\Pi_1$ at 1\,W, $\Pi_2$, $\Pi_3$ gaming, $\Pi_4$ social feed (with $v_1=0$),
$\Pi_5$ typical day. Other classes: illustrative mid values without device data (input ranges move each by about
a decade). Energy: nominal $E_n$.}
\label{tab:general}
\setlength{\tabcolsep}{4.5pt}
\begin{tabular}{@{}lrrrrrrr@{}}
\toprule
 & RP-63 & Tablet & Laptop & Watch & Earbuds & E-reader & Drone \\
\midrule
Energy (W\,h) & 18.2 & 30 & 55 & 1.1 & 0.2 & 5 & 40 \\
$\Pi_1$ ohmic loss & 0.0047 & 0.0085 & 0.0075 & 0.0078 & 0.011 & 0.0017 & 0.032 \\
$\Pi_2$ throttling & 1.19 & 0.89 & 1.8 & 1.2 & 0.13 & 0.17 & 0.30 \\
$\Pi_3$ thermal memory & 0.022 & 0.045 & 0.12 & 0.0077 & 0.0034 & 0.017 & 0.32 \\
$\Pi_4$ burst margin ($\soc=0.5$, 25\,$^\circ$C, new) & 0.11 & 0.092 & 0.15 & 0.31 & 0.11 & 0.13 & 0.32 \\
$\Pi_4$ ($\soc=0.2$, $-20\,^\circ$C, SOH 0.8) & 1.18 & 1.4 & 2.0 & 4.3 & 2.1 & 1.1 & $\infty$ \\
$\Pi_5$ usage spread & 0.175 & 0.24 & 0.18 & 0.16 & 0.039 & 0.079 & 0.034 \\
\bottomrule
\end{tabular}
\end{table}

RP-63 has a group above 1 only in gaming ($\Pi_2=1.19$) or when cold, aged and low ($\Pi_4=1.18$; 1.35 with RC
polarization). Drones have the largest $\Pi_1$ and $\Pi_3$, and, cold and aged, hover alone breaches
$V_{\mathrm{sys}}$ (pre-warm the pack). Watches ($\Pi_4=4.3$ cold) and e-readers are floor-dominated (TTE elasticity
to the floor about $-0.6$ and $-0.7$; an e-reader lasts about 48 days on 5\,W\,h): standby is their lever.
Earbuds, e-readers and drones keep $\Pi_5<0.1$. Every class reaches $\Pi_4\gtrsim1$ when cold, aged
and low, so O1 carries over.
